\item Pruebe la fórmula $A = \frac{1}{2}r^2\theta$ correspondiente al área de un sector circular de radio $r$ y ángulo central $\theta$. (\textit{Sugerencia:} Suponga que $0 < \theta < \pi/2$ y coloque el centro de la circunferencia en el origen de manera que tenga por ecuación $x^2 + y^2 = r^2$. Entonces $A$ es la suma del área del triángulo $POQ$ y el área de la región $PQR$ mostrados en la figura).
    
    \begin{center}
    \begin{tikzpicture}[scale=2.5]
        % Ejes
        \draw[->] (-0.2,0) -- (1.5,0) node[below] {$x$};
        \draw[->] (0,-0.2) -- (0,1.2) node[left] {$y$};
        \node[below left] at (0,0) {$0$};

        % Coordenadas
        \coordinate (O) at (0,0);
        \coordinate (P) at (45:1);
        \coordinate (Q) at ({cos(45)}, 0);
        \coordinate (R) at (1,0);

        % Arco y líneas
        \draw (0,1) arc (90:0:1);
        \draw (O) -- (P) node[above] {$P$};
        \draw (P) -- (Q);
        
        % Etiquetas
        \node[below] at (Q) {$Q$};
        \node[below] at (R) {$R$};

        % Ángulo recto en Q
        \draw ({cos(45)-0.05}, 0) -- ({cos(45)-0.05}, 0.05) -- ({cos(45)}, 0.05);

        % Ángulo theta
        \pic [draw, "$\theta$", angle eccentricity=1.5, angle radius=0.4cm] {angle = Q--O--P};
    \end{tikzpicture}
    \end{center}
